\section{Testing from VSCode}\label{sec:vscode}

This section is a step-by-step guide to running the plugin and checking its
certificates by hand in VSCode. It assumes nothing beyond a working opam
switch with Rocq 9 and the plugin's dependencies (\file{README.md},
``Installation'').

\paragraph{1. The editor and its extension.} Install VSCode and the
\emph{VsRocq} extension (the extensions panel, search ``VsRocq''). VsRocq
runs a Rocq server, \code{vsrocqtop}, which must be on your \code{PATH}:
\code{which vsrocqtop} in a terminal prints its path (for a local switch,
\file{README.md}, ``Using VSCode'', gives the setting that names it).

\paragraph{2. Build first, under a memory cap.} VsRocq does not build the
project: it reads \file{_RocqProject}, which maps the logical names
\code{MEBI} (\file{theories/}), \code{MEBI.Examples} and \code{MEBITests}
to dune's compiled files in \file{_build/}, and loads the plugin from
\file{_build/install/default/lib}. So, from the repository's root:
\begin{verbatim}
systemd-run --user --scope -p MemoryMax=1750M -p MemorySwapMax=0 \
  dune build -j 2
\end{verbatim}
(\code{systemd-run} caps the build's memory: a build that needs more is
stopped instead of pushing the machine into swap.) After any change to
the plugin, build again, then run the command \emph{VsRocq: Restart} (the
command palette, \code{Ctrl+Shift+P}).

\paragraph{3. Open the folder, then a file.} \emph{File, Open Folder}, the
repository's root: VsRocq finds \file{_RocqProject} there. Open
\file{test/rocq/LogTest.v}. VsRocq checks a file up to a point you choose:
\emph{Interpret to point} (\code{Alt+Right}) up to the cursor,
\emph{Step forward} (\code{Alt+Down}) one command, \emph{Interpret to end}
(\code{Alt+End}). The goals appear in the \emph{goal panel}, the plugin's
messages under the goals (and in \emph{View, Output}, channel
\emph{VsRocq}).

\paragraph{4. What to expect in \file{LogTest.v}.} Step through it command
by command:
\begin{itemize}
\item \code{Proof. weak_bisim. Qed.} The tactic decides the goal,
  builds the proof from the certificate and prints \code{weak_bisim:
  decided (12 pairs); the proof from Sylvan's certificate built in
  \ldots}; \code{Qed} has the kernel check it. With \code{Set MeBi Log
  Mode "keep"} before it, the log stays in
  \file{test/rocq/.LogTest.v.mebi/logged.log}, next to the file you opened
  (git ignores these directories).
\item \code{MeBi Check Log ".LogTest.v.mebi/logged.log".} The log's records
  become facts and the kernel checks them; the message gives the chunks,
  the facts, the nodes and the times, and the last chunk's constant,
  \code{mebi_log1_c1}.
\item \code{Print Assumptions mebi_log1_c1.} Must print \emph{Closed under
  the global context}: the facts use no axiom.
\end{itemize}

\paragraph{5. The theorems proved once.} In a scratch file (in
\file{_certify/}, which git ignores, so that nothing is left in the source
tree), check the transfer of \cref{sec:status}:
\begin{verbatim}
Require Import MEBI.DdGame.
Check cert_weak_sim.
Print Assumptions cert_weak_sim.
Print Assumptions cert_weak_bisim.
\end{verbatim}
\code{Check} shows the statement (the game on vectors, completeness for
the mover, soundness for the answerer, then \code{weak_sim}); each
\code{Print Assumptions} must print \emph{Closed under the global
context}.

\paragraph{6. Measuring from the editor.} The plugin's measurements go to
standard error, controlled by environment variables
(\code{MEBI_CERT_STATS=1}, \file{CLAUDE.md}). The Rocq server inherits
VSCode's environment, so close VSCode, then start it from a terminal with
them set, and capped (a running VSCode would be reused, without them):
\begin{verbatim}
MEBI_CERT_STATS=1 systemd-run --user --scope \
  -p MemoryMax=4G -p MemorySwapMax=0 code .
\end{verbatim}
The lines appear in the extension's channel of \emph{View, Output}. The cap
covers everything VSCode starts, the Rocq server included.

\paragraph{7. What not to step through.} The editor runs Rocq with no cap
of its own (unless started as above): step through small files only
(\file{test/rocq/}, the CCS and Proc examples). The CADP examples with
two and three processes are for \file{bench/prooflog.py}, which caps each
step.

\paragraph{8. Goals close.} A tactic ends its goal with the proof from
the certificate (\cref{sec:qed}): the test of a lemma \code{L} is
\code{Proof. weak_sim. Qed.} then \code{Print Assumptions L.}, which must
print \emph{Closed under the global context}.
